% !TEX root = /home/eb/.vscode/extensions/qftrules.programme-colle/tmp/Exercice.tex

%%%%%%%%%%%%%%%%%%%%%%%%%%%%%%%%%%%%%%%%%%%%%%%%%%%%%%%%%%%%%%%%%%%%%%%%%%%%
\begin{exo}[PCSI]{Adaptation d'impédance}
On étudie un réacteur à plasma modélisé par un circuit série $R_{\rm{p}}-C_{\rm{p}}$, d'impédance complexe $\underline{Z}_{\rm{p}}$. On cherche à diminuer au maximum la partie imaginaire de $\underline{Z}_{\rm{p}}$ et dans ce but, on réalise le circuit ci-dessous.
%\begin{figure}[h!]
%\centering
\begin{center}
\includegraphics[width=0.25\columnwidth]{adaptation_impedance.pdf}
\end{center}
%\end{figure}

\begin{quest}
\item Exprimer l'impédance complexe totale, notée $\underline{Z}$, du circuit ci-dessus.
\item En déduire la valeur de $C$ qui annule la partie imaginaire de l'impédance totale $\underline{Z}$.
\item La condition précédente étant réalisée, déterminer l'expression de l'impédance $\underline{Z}$ totale du dipôle.
\end{quest}
\end{exo}
%%%%%%%%%%%%%%%%%%%%%%%%%%%%%%%%%%%%%%%%%%%%%%%%%%%%%%%%%%%%%%%%%%%%%%%%%%%%

%%%%%%%%%%%%%%%%%%%%%%%%%%%%%%%%%%%
\begin{exo}[PC][1][colle]{Modulation d'amplitude}
On considère un signal modulé en amplitude de la forme $s(t)=k(A_0+s_i(t))s_p(t)$, où $A_0$ est une constante. Les signaux informatif $s_i(t) = A_i \cos(\omega_i t)$ et porteur $s_p(t) = A_p \cos(\omega_p t)$ sont des sinusoïdes de fréquences respectives $\omega_i$ et $\omega_p$ et d'amplitudes respectives $A_i$ et $A_p$.
\begin{quest}
    \item Déterminer le spectre en fréquence du signal modulé $s(t)$. La porteuse est-elle conservée ?
    \item Définir le taux de modulation $m$ et le calculer pour $A_i=\SI{1}{V}$ et $A_0=\SI{2}{V}$. Peut-on démoduler le signal par détection de crêtes ?
    \item Proposer un montage utilisant une diode et une circuit RC pour démoduler le signal. Justifier son fonctionnement. Expliquer l'effet du circuit RC sur le spectre du signal.
\end{quest}
\end{exo}
%%%%%%%%%%%%%%%%%%%%%%%%%%%%%%%%%%%

\begin{exo}[PC][1][colle]{Modulation de fréquence}
    On considère un signal modulé en fréquence de la forme $s(t)=A_0\cos(\omega_0 t + k s_i(t))$, où $A_0$ est une constante, $\omega_0=2\pi f_0$ est la pulsation de la porteuse et $k$ est la sensibilité du modulateur. Le signal informatif $s_i(t)$ est une sinusoïde de pulsation $\omega_i=2\pi f_i$ et d'amplitude $A_i$.
    \begin{quest}
        \item	Donner l'expression du signal modulé $s(t)$. Définir le taux de modulation $\beta$ et la variation de fréquence $\Delta f$ en fonction de $k$, $A_i$ et $f_i$. Donner l'allure du signal modulé $s(t)$ en fonction du temps.
        \item Donner l'allure du spectre en fréquence du signal modulé. Autour de quelle fréquence le spectre est-il centré ? Que vaut la largeur de bande $B$ du signal modulé (règle de Carson) ? 
    \end{quest}
\end{exo}
%%%%%%%%%%%%%%%%%%%%%%%%%%%%%%%%%%%

%%%%%%%%%%%%%%%%%%%%%%%%%%%%%%%%%%%%%%%%%%%%%%%%%%%%%%%%%%%%%%%
\begin{exo}[PCSI][2][colle]{Circuit bouchon}
Un circuit bouchon est constitué d'une bobine réelle $(r,L)$ en parallèle avec un condensateur de capacité $C$, alimenté par la tension sinusoïdale $e(t)=E_{0}\cos(\omega t)$. On travaillera avec les grandeurs complexes.
%\begin{figure}[h!]
%\centering
\begin{center}
\includegraphics[width=0.25\columnwidth]{bouchon.pdf}
\end{center}
%\end{figure}
\begin{quest}
\item Déterminer l'expression de l'impédance $\underline{Z}(\omega)$ du circuit bouchon.
\item Pour quelle pulsation $\omega_0$ a-t-on résonance ? Quel vaut alors la valeur $\udl{Z}(\omega_0)$ de l'impédance à la résonance ? 
\item Montrer que si le facteur de qualité vérifie $Q=\frac{L\omega_0}{r}\gg 1$, l'impédance à la résonance se met sous la forme $\udl{Z}(\omega_0) \simeq rQ^2$. Justifier alors la dénomination de circuit bouchon.
\item On pose $\omega - \omega_0 = \omega_0\, \veps$. On se place maintenant au voisinage de la résonance, de sorte que $\veps \ll 1$. Développer le rapport $\frac{\underline{Z}(\omega)}{\underline{Z}(\omega_0)}$ en fonction de $Q$ à l'ordre dominant en $\veps$.
\end{quest}
\end{exo}
%%%%%%%%%%%%%%%%%%%%%%%%%%%%%%%%%%%%%%%%%%%%%%%%%%%%%%%%%%%%%%%

%%%%%%%%%%%%%%%%%%%%%%%%%%%%%%%%%%%%%%%%%%%%%%%%%%%%%%%%%%%%%%%
\begin{exo}[PCSI]{Étude d'un quartz}
Les oscillateurs à quartz sont utilisés pour leur grande stabilité en fréquence. Ils sont notamment utilisés dans les montres électroniques, où ils sont taillés en diapason et réglés pour vibrer à une fréquence de \SI{32768}{Hz}.
Entre ses deux bornes $A$ et $B$, un quartz est équivalent au circuit électrique ci-dessous.
%\begin{figure}[h!]
%\centering
\begin{center}
    \includegraphics[width=0.4\columnwidth]{quartz.pdf}
\end{center}
%\end{figure}

\begin{quest}
\item Exprimer l'impédance complexe $\underline{Z}$ entre $A$ et $B$ en fonction de $L,C_{\rm{s}},C_{\rm{p}}$ et $\omega$. 
\item Mettre l'impédance sous la forme
\begin{equation}
    \underline{Z}=\frac{j(\omega_{\rm{s}}^2-\omega^2)}{C_{\rm{p}}\omega(\omega^2-\omega_{\rm{p}}^2)},
\end{equation} 
en exprimant les grandeurs $\omega_{\rm{s}}$ et $\omega_{\rm{p}}$  en fonction de $L$, $C_{\rm{s}}$ et $C_{\rm{p}}$. 
\item Représenter le graphe de la partie imaginaire de l'impédance en fonction de la pulsation. À quelle pulsation a-t-on résonance (transmission maximale de courant sans perte) ? Déterminer $L$ et $C_{\rm{p}}$ pour que cette résonance soit à \SI{32768}{Hz}.
\item Dans quel domaine le circuit a-t-il un comportement inductif ?  capacitif ? justifier.
\end{quest}
\end{exo}
%%%%%%%%%%%%%%%%%%%%%%%%%%%%%%%%%%%%%%%%%%%%%%%%%%%%%%%%%%%%%%%


%%%%%%%%%%%%%%%%%%%%%%%%%%%%%%%%%%%%%%%%%%%%%%%%%%%%%%%%%%
\begin{exo}[PCSI]{Filtre de Hartley}
%\begin{figure}[h!]
%\centering
\begin{center}
\includegraphics[width=0.45\columnwidth]{hartley.pdf}
\end{center}
%\end{figure}
\begin{quest}
\item Étudier qualitativement le comportement du filtre à basses et hautes fréquences.
\item Déterminer la fonction de transfert du filtre et la mettre sous forme canonique 
\begin{equation}
    \udl{H} = \frac{H_0 jx/Q}{1-x^2+jx/Q},
\end{equation}
avec $x=\omega/\omega_0$. Donner les expressions de $H_0$, du facteur de qualité $Q$ et de la pulsation propre $\omega_0$.
\item On pose maintenant $L_1=L_2=L=\SI{1}{mH}$, $C=\SI{0,1}{\micro F}$ et $R=\SI{10}{k\ohm}$. Tracer le diagramme de Bode complet du filtre (gain et phase),  d’abord en traçant les asymptotes, puis l’allure des courbes exactes.
\item On envoie en entrée un signal carré, dont le spectre est représenté ci-dessous.
\begin{fig}[label]
    \centering
    \includegraphics[width=0.5\columnwidth]{square_signal.jpg}
    %\capt{}
\end{fig}
Représenter l'allure du spectre du signal en sortie.
\end{quest}
\end{exo}
%%%%%%%%%%%%%%%%%%%%%%%%%%%%%%%%%%%%%%%%%%%%%%%%%%%%%%%%%%

%%%%%%%%%%%%%%%%%%%%%%%%%%%%%%%%%%%%%%%%%%%%%%%%%%%%%%%%%%
\begin{exo}[PCSI]{Filtre de Wien}
    \Notions{\item Électrocinétique \task Filtre passif \task Condensateur \task Diagramme de Bode}
    On considère le filtre passif suivant. Les résistances $R$ sont identiques, de même que les capacités $C$. La tension d’entrée est notée $e$ et la tension de sortie $s$.
%\begin{figure}[h!]
%\centering
\begin{center}
\includegraphics[width=0.45\columnwidth]{filtre_wien.pdf}
\end{center}
%\end{figure}
\begin{quest}
\item Étudier qualitativement le comportement du filtre à basses et hautes fréquences.
\sol{
  \begin{liste}
    \item
    Lorsque $\omega \rightarrow 0$, les condensateurs se comportent comme des interrupteurs ouverts. Ainsi aucun courant ne circule, ni dans la branche d'entrée, ni dans la résistance en parallèle, et donc $s=0$.
    \item Lorsque $\omega \rightarrow \infty$, les condensateurs se comportent comme des fils. Or $s$ correspond alors à la tension aux bornes d'un fil, et donc $s=0$.
  \end{liste}

Le filtre est donc un \fbox{passe-bande}.
}
\item Déterminer la fonction de transfert du filtre et la mettre sous forme canonique, comme 
\begin{equation}
    \udl{H} = \frac{H_0}{1+jQ(x-\frac{1}{x})},
\end{equation}
avec $x=\omega/\omega_0$. Donner les expressions de $H_0$, du facteur de qualité $Q$ et de la pulsation propre $\omega_0$.
\sol{
    Notons $Z_s = R+\frac{1}{jC\omega}$ la résistance de l'association de $R$ et $C$ en série, et $Z_p = \frac{R}{1+jRC\omega}$ la résistance de l'association de $R$ et $C$ en parallèle. Un pont diviseur de tension donne
$$
\udl{H} = \frac{\udl{s}_f}{\udl{e}_f} = \frac{Z_p}{Z_s+Z_p} = \frac{R}{R+(\frac{1}{jC\omega}+R)(1+{jRC\omega})} = \frac{R}{3R+jR^2C\omega - j\frac{\omega}{C}},
$$
d'où l'on tire 
% \eqbox{
$$
\udl{H} = \frac{1}{3}\frac{1}{1+j\frac{1}{3}\left(\frac{\omega}{1/RC}-\frac{1/RC}{\omega}\right)}.
$$
% }
On peut mettre la fonction de transfert sous la forme canonique demandée en posant
\eqbox{
H_0 = \frac{1}{3}, \quad Q = \frac{1}{3}, \quad \omega_0 = \frac{1}{RC}.
}
On trouve bien l'expression demandée avec $x=\omega/\omega_0$.
}
\item On donne $C=\SI{10}{nF}$ et $R=\SI{10}{k\ohm}$. Calculer $H_0$, $\omega_0$ et $Q$. Tracer le diagramme de Bode complet du filtre (gain et phase), d’abord en traçant les asymptotes, puis l’allure des courbes exactes.
\sol{
    \begin{enumerate}
        \item \AN~: $H_0 \simeq \num{0.33}, \quad Q \simeq \num{0.33}, \quad \omega_0 =\SI{1.0e4}{rad.s^{-1}}$.
        \item On rappelle la définition du gain et de la phase :
        % $$
        \eqbox{
            G_{dB} = 20 \log |\udl{H}|, \quad \varphi = \arg(\udl{H}).
            }
        % $$
        Les asymptotes s'obtiennent dans les limites haute et basse fréquences. On trouve 
        \begin{itemize}
            \item Lorsque $x \rightarrow 0$, on a $\udl{H} \sim j\frac{H_0}{Q}x$, donc $G_{dB} \simeq +20 \log(x)$ et $\varphi \simeq +\frac{\pi}{2}$.
            \item Lorsque $x \rightarrow \infty$, on a $\udl{H} \sim -j \frac{H_0}{Qx}$, donc $G_{dB} \simeq -20 \log(x)$ et $\varphi \simeq -\frac{\pi}{2}$.
            \item Lorsque $x=1$, on a $\udl{H} = H_0$, donc $G_{dB} = 20 \log(H_0) = - 20 \log 3 \simeq -\num{9.5}$, et $\varphi = 0$.
        \end{itemize}
        Les courbes exactes sont tracées ci-dessous. 
        \begin{fig}[label]
            \centering
            \includegraphics[width=0.7\columnwidth]{GdB_Wien}
            %\capt{}
        \end{fig}
        \begin{fig}[label]
            \centering
            \includegraphics[width=0.7\columnwidth]{phase_Wien}
            %\capt{}
        \end{fig}
    \end{enumerate}
}
\item On envoie en entrée un signal carré, dont le spectre est représenté ci-dessous. L'axe des abscisses correspond à la variable $x=\omega/\omega_0$.
\begin{fig}[label]
    \centering
    \includegraphics[width=0.5\columnwidth]{square_signal.jpg}
    %\capt{}
\end{fig}
Représenter l'allure du spectre du signal en sortie.
\sol{
La fréquence fondamentale du signal en entrée correspond à la fréquence propre du filtre de Wien. On voit que le signal contient ensuite toutes les harmoniques impaires. Si $A_k$ désigne l'amplitude du k-ème harmonique, de pulsation adimensionnée $x_k=2k+1$, l'amplitude en sortie doit valoir
$$
B_k = {\udl{H}(x_k)} A_k \simeq \frac{A_k}{2k+1}.
$$

}
\end{quest}
\end{exo}
%%%%%%%%%%%%%%%%%%%%%%%%%%%%%%%%%%%%%%%%%%%%%%%%%%%%%%%%%%

%%%%%%%%%%%%%%%%%%%%%%%%%%%%%%%%%%%%%%%%%%%%
\begin{exo}[MPSI]{Filtres RC en cascade}
Le filtre suivant est réalisé en associant deux filtres RC en cascade. Les deux filtres sont reliés par un montage suiveur (non discuté ici) qui permet d'obtenir la tension $V_{\rm i}$ à l'entrée du filtre de droite sans tirer du courant du filtre de gauche. On envoie une tension sinusoïdale en entrée, notée $V_{ \rm e}$. On récupère la tension $V_{\rm s}$ en sortie du filtre.

%%% FIGURE %%%
%\begin{figure}[h!]
%\centering
\begin{center}
\shorthandoff{:!}
\begin{tikzpicture}

\draw (0,0) coordinate (O) to[C=$C$] ++ (1,0) coordinate (A);
\draw[white] (O) -- ++ (0,-1.5) -- ++ (2,0) coordinate (B);

\draw (A) -- ++ (1.5,0) coordinate (D) to[R=$R$] ++ (0,-1.5) -- (B);

\draw (O) to[battery1] ++(0,-1.5) coordinate (R) node[ground]{};

\draw (R) -- (B);

\draw[<-,>=stealth]  ($(D)+(0.75,0)$) node[below right]{$V_{\rm i}$} -- ++ (0,-1.5) coordinate (G);
\draw[<-,>=stealth]  ($(O)+(-0.5,0)$) node[left]{$V_{\rm e}$} -- ++ (0,-1.5) coordinate (H);

\draw (D) to[short, i=${i=0}$] ++ (3,0) coordinate (E);

\draw (E) -- ++ (0,0.5) -- ++ (1.5,0) node[midway,below=5] {suiveur} -- ++ (0,-1) -- ++ (-1.5,0) -- cycle;

\draw (7,0) -- ++ (0.5,0) to[R=$R$] ++ (2,0) coordinate (Dp) to[C=$C$] ++ (0,-1.5) node[ground] {};

\draw[<-,>=stealth]  ($(Dp)+(1,0)$) node[right]{$V_{\rm s}$} -- ++ (0,-1.5) coordinate (H);



\end{tikzpicture}
\end{center}
%\caption{Les fentes d'Young $S_1$ et $S_2$.}
%\label{fig:}
%\end{figure}
%%%%%%%%%%

\begin{quest}
\item Sans calcul, étudier les comportements haute et basse fréquence ($\omega=\infty$ et $\omega=0$) de ce filtre.

\sol{
Condensateur = fil pour $\omega=\infty$ donc $V_s=0$. Condensateur = circuit ouvert pour  $\omega=0$ donc $V_s=0$.
}


\item Déterminer la fonction de transfert $H_1$ du filtre de gauche (en entrée du suiveur). Quelle est la nature de ce filtre ?

\sol{
$H_1 = \dfrac{R}{R+1/jC\omega} = \dfrac{jRC\omega}{1+jRC\omega}$. Passe-haut.
}

\item Déterminer la fonction de transfert $H_2$ du filtre de droite (en sortie du suiveur). Quelle est la nature de ce filtre ?

\sol{
$H_2 = \dfrac{1/jC\omega}{R+1/jC\omega} = \dfrac{1}{1+jRC\omega}$. Passe-bas.
}

\item En déduire la fonction de transfert $H$ du filtre complet. Définir une pulsation caractéristique du circuit, notée $\omega_0$. Représenter le diagramme de Bode en gain et en phase du filtre complet. Comment pourraît-on appeler ce filtre ?

\sol{
$H = H_1H_2$
}

\end{quest}
\end{exo}
%%%%%%%%%%%%%%%%%%%%%%%%%%%%%%%%%%%%%%%%%

%%%%%%%%%%%%%%%%%%%%%%%%%%%%%%%%%%%%%%%%%%%%%%%%%%%%%%%%%%%%%%
\begin{exo}[MPSI]{Circuit RL en régime forcé}

On souhaite étudier la réponse du circuit suivant à une entrée sinusoïdale $V_e(t)=E\cos(\omega t)$ pour $t>0$. La tension d'entrée est nulle pour $t<0$. On se propose d'étudier le circuit en notations réelle et complexe afin de comparer les deux approches.

%%% FIGURE %%%
%\begin{figure}[h!]
%\centering
\begin{center}
\shorthandoff{:!}
\begin{tikzpicture}

\draw (0,0) coordinate (O) to[L=$L$] ++ (3,0) coordinate (A);
\draw[white] (O) -- ++ (0,-1.5) -- ++ (3,0) coordinate (B);
\draw (A) to[R=$3R$] (B);

\draw (A) to[R=$R$] ++ (2,0) coordinate (D) to[R=$R$] ++ (0,-1.5) -- (B);

\draw (O) to[battery1] ++(0,-1.5) coordinate (R) node[ground]{};

\draw (R) -- (B);

\draw[<-,>=stealth]  ($(A)-(0.5,0)$) node[below left]{$V_{\rm s}$} -- ++ (0,-1.5) coordinate (G);
\draw[<-,>=stealth]  ($(O)+(-0.5,0)$) node[left]{$V_{\rm e}$} -- ++ (0,-1.5) coordinate (H);


\end{tikzpicture}
\end{center}
%\caption{}
%\label{fig:}
%\end{figure}
%%%%%%%%%%

\begin{quest}
\item Déterminer la résistance équivalente $R_{\rm eq}$ à l'association des trois résistances du bloc de droite du circuit. Représenter le schéma équivalent du circuit. On travaillera avec la grandeur $R_{\rm eq}$ par la suite.

\sol{
$R_{\rm eq} = 2R \parallel 3R = \dfrac{6}{5} R$.
}



\item  Équation d'évolution.


\begin{enumerate}[label=\alph{*})]

\item À partir des lois de Kirchhoff, établir l'équation différentielle qui lie la tension de sortie $V_{\rm s}(t)$ à la tension d'entrée $V_{\rm e}(t)$. Définir le temps caractéristique $\tau$ qui intervient dans cette équation.

\sol{
$V_e = V_s + L\d i/\d t = V_s + \dfrac{L}{R_{\rm eq}} \d V_s/\d t$.
}

\item À partir des lois de Kirchhoff en notation complexe, déterminer l'amplitude complexe $\underline{V_s}$ de la tension de sortie en fonction de l'amplitude de la tension d'entrée $E$. Peut-on décrire un régime transitoire avec cette équation ?

\end{enumerate}


\item Nous étudions le régime transitoire en notation réelle. Que vaut la période $T$ du signal d'entrée en fonction de $\omega$ ? On suppose que le temps caractéristique $\tau$ est très petit devant la période~$T$, de sorte à ce qu'on puisse supposer $V_e(t)$ constante, égale à $E$, sur une durée de quelques $\tau$. Résoudre alors l'équation différentielle pour obtenir $V_s(t)$ en régime transitoire.

\sol{
$i$ et $V_s$ continues. $V_s(t) = E(1-e^{-t/\tau})$ avec $\tau = L/R_{\rm eq}$.
}


\item Régime forcé.
\begin{enumerate}[label=\alph{*})]
\item On étudie le régime forcé en notation complexe. À partir de la question 2)b), déterminer le module $A$ et la phase $\phi$ de la tension de sortie complexe $\underline{V_s}$.


\item On souhaite maintenant étudier le régime forcé en notation réelle. Chercher une solution à l'équation différentielle de la question 2)a) sous la forme $V_s(t) = A \cos(\omega t +\phi)$. Trouver les expressions de $A$ et $\phi$. Comparer les deux méthodes.

\end{enumerate}

\end{quest}
\end{exo}
%%%%%%%%%%%%%%%%%%%%%%%%%%%%%%%%%%%%%%%%%%

%%%%%%%%%%%%%%%%%%%%%%%%%%%%%%%%%%%%%%%%%%
\begin{exo}[MPSI]{Filtres passe-bas en cascade}

Le filtre suivant est réalisé en associant deux filtres RL et RC en cascade. Les deux filtres sont reliés par un montage suiveur (non discuté ici) qui permet d'obtenir la tension $V_{\rm i}$ à l'entrée du filtre de droite sans tirer du courant du filtre de gauche. On envoie une tension sinusoïdale en entrée, notée~$V_{ \rm e}$. On récupère la tension $V_{\rm s}$ en sortie du filtre.

%%% FIGURE %%%
%\begin{figure}[h!]
%\centering
\begin{center}
\input{RL.pgf}
\end{center}
%\caption{Les fentes d'Young $S_1$ et $S_2$.}
%\label{fig:}
%\end{figure}
%%%%%%%%%%

\begin{quest}
\item Sans calcul, étudier les comportements haute et basse fréquence ($\omega=\infty$ et $\omega=0$) de ce filtre.

\sol{
Condensateur = fil et bobine = ouvert pour $\omega=\infty$ donc $V_s=0$. Condensateur = circuit ouvert pour $\omega=0$ donc $V_s=V_e$.
}


\item Déterminer la fonction de transfert $H_1$ du filtre de gauche (en entrée du suiveur). Quelle est la nature de ce filtre ? Déterminer la pulsation caractéristique de ce filtre, notée $\omega_1$.

\sol{
$H_1 = \dfrac{R}{R+jL\omega} = \dfrac{1}{1+j\omega/\omega_1}$ avec $\omega_1 = R/L$. Passe-bas.
}

\item Déterminer la fonction de transfert $H_2$ du filtre de droite (en sortie du suiveur). Quelle est la nature de ce filtre ?  Déterminer la pulsation caractéristique de ce filtre, notée $\omega_2$.

\sol{
$H_2 = \dfrac{1/jC\omega}{R+1/jC\omega} = \dfrac{1}{1+j\omega/\omega_2}$ avec $\omega_2 = R/L$. Passe-bas.
}

\item En déduire la fonction de transfert $H$ du filtre complet. Tracer le diagramme de Bode en gain et en phase du filtre complet pour $L=2\,$mH, $C=100\,$nF et $R=2\,$k$\Omega$.

\sol{
$H = H_1H_2$
}

\end{quest}
\end{exo}
%%%%%%%%%%%%%%%%%%%%%%%%%%%%%%%%%%%%%%%%%%%%%
